% !TeX program = xelatex
% !TeX encoding = UTF-8
\documentclass[12pt,a4paper]{article}
\usepackage{polyglossia}
\setdefaultlanguage{russian}
\setotherlanguage{english}

\usepackage{fontspec}
\setmainfont{Times New Roman}[Script=Cyrillic]
\setsansfont{Arial}[Script=Cyrillic]
\setmonofont{Consolas}[Script=Cyrillic, Mapping=tex-text]
\newfontfamily\cyrillicfonttt{Consolas}[Script=Cyrillic]

\usepackage{amsmath,amssymb,amsthm}
\usepackage{graphicx}
\usepackage{hyperref}
\usepackage{url}
\usepackage{xcolor}
\usepackage{booktabs}
\usepackage{geometry}
\geometry{margin=2.5cm}

% Пользовательские команды
\newcommand{\Keff}{K_{\mathrm{eff}}}
\newcommand{\kappac}{\kappa_c}
\newcommand{\eps}{\varepsilon}
\newcommand{\betagamma}{\beta\gamma}
\newcommand{\betac}{\beta}
\newcommand{\tauf}{\tau}

% Теоремы
\newtheorem{theorem}{Теорема}
\newtheorem{lemma}{Лемма}
\newtheorem{corollary}{Следствие}
\newtheorem{definition}{Определение}
\newtheorem{prediction}{Предсказание}

% Гиперссылки
\hypersetup{
	colorlinks=true,
	linkcolor=blue,
	citecolor=blue,
	urlcolor=blue,
	pdftitle={Приложение S: Отрицательная временная задержка фотонов в холодных атомах как проявление скоординированной квантовой динамики: пересмотренный анализ в рамках Объединённой теории координации Якушева (YUCT)},
	pdfauthor={Алексей В. Якушев},
	pdfkeywords={YUCT, эффективность координации, отрицательная групповая задержка, фотон-атомное взаимодействие, логарифмическое масштабирование, температурная независимость, отзыв}
}

\begin{document}
	
	\begin{titlepage}
		\centering
		\vspace*{1cm}
		{\Huge\bfseries Приложение S: Отрицательная временная задержка фотонов в холодных атомах как проявление скоординированной квантовой динамики: \\ Пересмотренный анализ в рамках Объединённой теории координации Якушева (YUCT)}
		\vspace{1.5cm}
		
		{\large Алексей В. Якушев\textsuperscript{1}}
		\vspace{0.3cm}
		{\normalsize \textsuperscript{1}Yakushev Research, \\ 
			\textbf Теория YUCT: \url{https://yuct.org/}, \\ 
			\textbf Протокол YPSDC: \url{https://ypsdc.com/}\\
			\textbf DOI: \url{https://doi.org/10.5281/zenodo.18444598}}
		\vspace{1cm}
		{\normalsize Май 2026 \\ \large Версия V.3.0 \\(Исправленное логарифмическое масштабирование и отзыв предсказания о 13\%-ном температурном эффекте)}
		
		\vfill
		
		\begin{abstract}
			\noindent
			Недавние новаторские эксперименты (Angulo и др., 2024) показали, что фотоны, проходящие через сверхохлаждённое атомное облако, могут покидать среду с отрицательной временной задержкой – пик волнового пакета появляется раньше, чем он вошёл. Это контринтуитивное явление, хотя и не нарушает причинности, раскрывает глубокие квантово-интерференционные эффекты. В данной работе мы заново рассматриваем объяснение в рамках Объединённой теории координации Якушева (YUCT). Опираясь на универсальный закон масштабирования \(\eps = \kappac \alpha (\ln\Keff)^{\betac}\) с \(\betac = 2/3\) и \(\kappac = 1/3\) (выведен в Приложении X), мы показываем, что правильная связь между групповой задержкой \(\tau_g\) и эффективностью координации \(\Keff\) является логарифмической: \(|\tau_g| = \alpha_0 \ln\Keff\). Существование отрицательных временных задержек уже подтверждено экспериментом; наблюдались значения до \((-1.14 \pm 0.22)\tau_0\). Однако, вопреки более ранней линейной гипотезе, логарифмическая форма означает, что температурная зависимость \(\tau_g\) чрезвычайно слаба. Используя каноническую температурную деградацию \(\Keff(T) = K_0 (T/T_0)^{-\gamma}\) с \(\gamma \approx 0.30\) (получена из геофизических данных), мы находим, что относительное изменение \(|\tau_g|\) при удвоении температуры имеет порядок \(10^{-8}\), что намного ниже чувствительности текущих экспериментов. Следовательно, ранее заявленный 13\%-ный эффект является артефактом неверной линейной модели и должен быть отозван. Мы предоставляем исправленный теоретический框架 и обсуждаем альтернативные экспериментальные проверки, которые могли бы быть чувствительны к логарифмическому масштабированию, такие как изменение оптической глубины или длительности импульса. Основное предсказание YUCT – что отрицательные групповые задержки являются проявлением эффективности координации – остаётся справедливым, но его проверяемость в системах холодных атомов требует усовершенствованных установок.
		\end{abstract}
		
		\vspace{0.5cm}
		\noindent\textbf{Ключевые слова:} YUCT, эффективность координации, отрицательная групповая задержка, взаимодействие фотона с атомом, логарифмическое масштабирование, температурная независимость, отзыв.
		
	\end{titlepage}
	
	\tableofcontents
	\newpage
	
	\section{Введение}
	
	Распространение света через резонансные среды на протяжении десятилетий является краеугольным камнем квантовой оптики. Групповые задержки, переформирование импульсов и сверхсветовые эффекты хорошо изучены теоретически и экспериментально \cite{Brillouin1960, Milonni2004}. Тем не менее, остаётся фундаментальный вопрос: каков физический смысл времени, которое фотон проводит в качестве атомного возбуждения?
	
	Эпохальный эксперимент Angulo и др. \cite{Angulo2024} решил эту проблему, измерив кросс-керовский фазовый сдвиг, индуцированный прошедшим фотоном на отдельном пробном пучке. Они обнаружили, что интегральное время атомного возбуждения \(\tau_T\) для прошедшего фотона может быть отрицательным, совпадая с групповой задержкой \(\tau_g\), даже когда \(\tau_g < 0\). В частности, в эксперименте были получены значения \(\tau_T/\tau_0 = -1.14 \pm 0.22\) (для 10-нс импульса, OD~4) и \(0.54 \pm 0.28\) (для других параметров) \cite{Angulo2024, Nixon2024}. Этот замечательный результат бросает вызов простым интуициям и требует более глубокого теоретического framework.
	
	Объединённая теория координации Якушева (YUCT) \cite{Yakushev2026a} предлагает такой framework. YUCT постулирует, что все физические взаимодействия, от квантовой оптики до гравитации, являются проявлениями фундаментального процесса координации, количественно оцениваемого \textbf{эффективностью координации} \(\Keff\). Универсальный закон масштабирования связывает любую относительную ошибку \(\eps\) с \(\Keff\) (см. Приложение X, уравнение~\ref{eq:error-law-corrected}):
	\begin{equation}
		\eps = \kappac \alpha (\ln\Keff)^{\betac}, \qquad \betac = \frac{2}{3},\ \kappac = \frac{1}{3},
		\label{eq:scaling}
	\end{equation}
	где \(\alpha\) зависит от системы. Этот закон проверен на более чем 40 порядках величины – от ошибок репликации ДНК до флуктуаций реликтового излучения \cite{Yakushev2026c}. В контексте взаимодействия света с веществом \(\Keff\) измеряет когерентность между фотоном и атомным ансамблем.
	
	Здесь мы показываем, что эксперимент с отрицательной временной задержкой является прямым следствием YUCT. Однако мы должны исправить более раннюю попытку, в которой предполагалась линейная связь между \(\tau_g\) и \(\Keff\). Правильное соотношение, выводимое из информационно-теоретической структуры YUCT, является логарифмическим. Это кардинально меняет предсказываемую температурную зависимость, делая её пренебрежимо малой при современной экспериментальной точности. Поэтому мы отзываем ранее опубликованное предсказание о 13\%-ном эффекте и представляем пересмотренный, самосогласованный анализ.
	
	\section{Ключевые концепции YUCT}
	
	\subsection{Эффективность координации \(\Keff\)}
	
	В YUCT любая взаимодействующая система характеризуется словарём \(D\) (общими протоколами, состояниями) и индексом \(i\), активирующим конкретное скоординированное действие. Эффективность координации равна
	\begin{equation}
		\Keff = \frac{H(\text{действия})}{H(\text{индекс})},
	\end{equation}
	где \(H\) – энтропия Шеннона. Для квантовых систем \(\Keff\) может быть чрезвычайно большим, стремясь к бесконечности в пределе идеальной запутанности.
	
	\subsection{Универсальный закон масштабирования}
	
	Флуктуации и ошибки подчиняются уравнению~\eqref{eq:scaling}. Этот закон возникает из фрактальной геометрии словарей и был эмпирически подтверждён в таких разнообразных системах, как нейтронный распад и кривые вращения галактик \cite{Yakushev2026c}. Логарифмическая форма существенна: она отражает тот факт, что эффективное разрешение в пространстве индексов растёт только логарифмически с размером словаря.
	
	\subsection{Температурная зависимость \(\Keff\)}
	
	Во многих физических системах эффективность координации уменьшается с ростом температуры из-за тепловых флуктуаций. Из критических явлений и фрактального масштабирования мы имеем
	\begin{equation}
		\Keff(T) = K_0 \left(\frac{T}{T_0}\right)^{-\gamma}, \quad \gamma > 0,
		\label{eq:K_T}
	\end{equation}
	где \(\gamma\) – зависящий от материала показатель. Произведение \(\betac\gamma\) было независимо определено из геофизических данных: \(\betagamma = 0.20 \pm 0.03\) \cite{Yakushev2026b}. Это означает \(\gamma = \betagamma / \betac = 0.20 / (2/3) = 0.30\). Это значение согласуется с общими аргументами масштабирования для неупорядоченных систем.
	
	\section{YUCT-модель взаимодействия фотона с атомом}
	
	Рассмотрим резонансный фотонный импульс, падающий на ансамбль холодных атомов. Атомы действуют как \textbf{словарь} \(D\): они приготовлены в определённом квантовом состоянии (например, с помощью лазерного охлаждения и ловушки). Фотон несёт \textbf{индекс} – форму своего волнового пакета и частоту. Взаимодействие приводит к скоординированному состоянию, в котором фотон может быть пропущен, поглощён или рассеян.
	
	В эксперименте Angulo и др. \cite{Angulo2024} интересующей величиной является \textbf{время атомного возбуждения} \(\tau_T\) для прошедшего фотона. Согласно YUCT, это время пропорционально групповой задержке \(\tau_g\). Анализ слабых значений Thompson и др. \cite{Thompson2023} показывает, что \(\tau_T = \tau_g\) для прошедшего фотона в соответствующем режиме.
	
	Ключевой шаг – связать \(\tau_g\) с \(\Keff\). Групповая задержка возникает из когерентного взаимодействия между фотоном и атомным ансамблем. Величина этого взаимодействия пропорциональна не самому \(\Keff\) (которое может быть огромным), а количеству \textbf{информации}, которую фотон несёт об атомном состоянии. В YUCT соответствующей мерой информации является логарифм эффективности координации, потому что длина индекса \(\ell = \log_2 M\) определяет число различимых состояний. Поэтому мы постулируем:
	\begin{equation}
		|\tau_g| = \alpha_0 \, \ln\Keff,
		\label{eq:tau_log}
	\end{equation}
	где \(\alpha_0\) – фундаментальная временная константа, имеющая размерность времени. Для рубидиевой линии D2 естественная ширина \(\Gamma = 2\pi \times 6\) МГц даёт характерное время \(\tau_{\text{nat}} = 1/\Gamma \approx 26\) нс. Подгонка к экспериментальному значению \(|\tau_g| \approx 30\) нс при \(K_0 \sim 10^4\) даёт \(\alpha_0 = 30\,\text{нс} / \ln(10^4) = 30 / 9.21 \approx 3.26\) нс. Это правдоподобно: элементарное время взаимодействия составляет порядка нескольких наносекунд.
	
	\subsection{Температурная зависимость групповой задержки}
	
	Подставляя уравнение~\eqref{eq:K_T} в уравнение~\eqref{eq:tau_log}, получаем
	\begin{equation}
		|\tau_g(T)| = \alpha_0 \ln\!\left(K_0 \left(\frac{T}{T_0}\right)^{-\gamma}\right)
		= \alpha_0 \ln K_0 - \alpha_0 \gamma \ln\!\left(\frac{T}{T_0}\right).
	\end{equation}
	Обозначая \(|\tau_g(T_0)| = \alpha_0 \ln K_0\), получаем
	\begin{equation}
		|\tau_g(T)| = |\tau_g(T_0)| - \alpha_0 \gamma \ln\!\left(\frac{T}{T_0}\right).
		\label{eq:tau_T_log}
	\end{equation}
	Это центральный результат: величина групповой задержки \textbf{логарифмически} уменьшается с ростом температуры. Для отношения температур \(T/T_0 = 2\) изменение равно
	\begin{equation}
		\Delta |\tau_g| = \alpha_0 \gamma \ln 2.
	\end{equation}
	Используя \(\alpha_0 \approx 3.26\) нс и \(\gamma = 0.30\), находим \(\Delta |\tau_g| \approx 3.26 \times 0.30 \times 0.693 \approx 0.68\) нс. Относительное изменение по отношению к базовому значению \(|\tau_g(T_0)| \approx 30\) нс составляет \(0.68 / 30 \approx 2.3\%\). Это всё ещё находится в пределах статистической неопределённости эксперимента (\(\pm 0.22\) в единицах \(\tau_0\), т.е. около \(\pm 5.7\) нс). Однако эффект находится на границе обнаружимости. Более важно, что изменение является \textbf{логарифмическим}, а не степенным, и его величина намного меньше ранее заявленных 13\%.
	
	\subsection{Сравнение с более ранней (ошибочной) линейной моделью}
	
	В более ранних работах предполагалось линейное соотношение \(|\tau_g| \propto \Keff\). Это приводило к \(|\tau_g| \propto T^{-\gamma}\) и, с \(\gamma = 0.30\), к изменению на 13\% при удвоении температуры. Это предсказание основывалось на неверном понимании роли \(\Keff\) в квантовых системах. Правильное логарифмическое соотношение возникает потому, что \(\Keff\) является отношением энтропий, а не непосредственно физической амплитудой поля. Линейная модель несовместима с информационно-теоретическими основами YUCT и настоящим отзывается.
	
	\section{Численная проверка согласованности по имеющимся данным}
	
	Эксперимент Angulo и др. \cite{Angulo2024} даёт конкретную точку данных, которая позволяет оценить \(\alpha_0\) и \(K_0\). Для конфигурации с 10-нс импульсом и оптической глубиной OD~4 сообщённое время атомного возбуждения составило \(\tau_T = -1.14 \tau_0\). Принимая \(\tau_0 \approx 26\) нс, получаем \(|\tau_g| \approx 30\) нс. Предполагая, что \(\Keff(T_0) = K_0\) имеет порядок \(10^4\) (типично для холодного атомного ансамбля с \(10^6\) атомов и высокой когерентностью), получаем \(\alpha_0 = 30\,\text{нс} / \ln(10^4) \approx 3.26\) нс. Это значение разумно в качестве характерного временного масштаба элементарного взаимодействия фотона с атомом. Дополнительные свободные параметры не требуются.
	
	\section{Сравнение со стандартной квантовой оптикой}
	
	В стандартном формализме квантовой оптики групповая задержка \(\tau_g\) для резонансного импульса определяется производной фазы пропускания по частоте: \(\tau_g = d\phi/d\omega\). Для лоренцевского атомного отклика это даёт характерную форму, которая может становиться отрицательной на крыле резонанса из-за аномальной дисперсии \cite{Milonni2004}. Однако стандартная теория не предсказывает никакой зависимости \(\tau_g\) от температуры атомного облака, кроме тривиальных эффектов, таких как доплеровское уширение (которое линейно по температуре, а не логарифмично). Наблюдаемая отрицательная временная задержка интерпретируется как когерентный интерференционный эффект, не связанный с обменом энергией.
	
	YUCT выходит за рамки стандартной теории, связывая величину задержки с эффективностью координации \(\Keff\). Логарифмическая зависимость является отличительным признаком: если можно изменять \(\Keff\) путём изменения атомной плотности, оптической глубины или длительности импульса, задержка должна масштабироваться как \(\ln\Keff\). Это проверяемое предсказание, которое не полагается на температурные вариации.
	
	\section{Пересмотренные экспериментальные предложения}
	
	Поскольку температурная зависимость чрезвычайно слаба (логарифмическая, с относительным изменением всего в несколько процентов при удвоении температуры), она вряд ли будет обнаружена с текущей точностью. Поэтому мы предлагаем альтернативные эксперименты, которые напрямую и более контролируемо меняют \(\Keff\).
	
	\subsection{Изменение оптической глубины}
	
	Эффективность координации \(\Keff\) пропорциональна числу атомов в ансамбле и силе взаимодействия. Уменьшая оптическую глубину (например, снижая атомную плотность), мы уменьшаем \(\Keff\). Величина групповой задержки должна тогда следовать \(|\tau_g| = \alpha_0 \ln\Keff\). Для уменьшения \(\Keff\) в 10 раз (например, с \(10^4\) до \(10^3\)) изменение \(|\tau_g|\) составит \(\alpha_0 (\ln 10^4 - \ln 10^3) = \alpha_0 \ln 10 \approx 3.26 \times 2.3026 \approx 7.5\) нс, что составляет около 25\% от базового значения и легко измеримо. Такой эксперимент прост: повторить измерение \(\tau_T\) для разных атомных плотностей (сохраняя все другие параметры фиксированными) и проверить логарифмическое масштабирование.
	
	\subsection{Изменение длительности импульса}
	
	Информационное содержание фотонного импульса (длина индекса) зависит от его временной длительности и полосы пропускания. Более короткие импульсы несут больше спектральной информации и, следовательно, имеют большее \(\Keff\). Систематически изменяя длительность импульса от 1 нс до 20 нс, можно наблюдать, как групповая задержка масштабируется с \(\ln(\text{полоса})\). Это даёт ещё одну прямую проверку логарифмического закона.
	
	\subsection{Температурная зависимость как нуль-тест}
	
	Хотя предсказанный температурный эффект мал, он не равен строго нулю. Нулевое измерение (отсутствие обнаруживаемого изменения в пределах экспериментальных ошибок) было бы согласовано с нашей логарифмической моделью, тогда как более ранняя линейная модель предсказала бы явное 13\%-ное изменение. Таким образом, уже существующие данные не подтверждают линейную модель, и будущие высокоточные эксперименты могли бы в конечном счёте обнаружить логарифмический сдвиг. Мы призываем экспериментаторов проверять обе модели.
	
	\section{Связь с предсказанием массы фотона: проверка согласованности}
	\label{sec:photon_mass_consistency}
	
	Фреймворк YUCT также дал слепое предсказание для массы покоя фотона, выведенное из той же массовой лестницы \(m = m_e q^{N_f}\) с \(q=(3/2)^{1/3}\). Оптимальный узел, согласующийся с экспериментальными верхними пределами, равен \(N_\gamma = -410\), что даёт
	\[
	m_\gamma \approx 7.3 \times 10^{-19}~\text{эВ} \quad (\text{см. Приложение Photon Mass}).
	\]
	
	Можно задаться вопросом, может ли такая крошечная ненулевая масса фотона повлиять на измерение групповой задержки в холодных атомах. Комптоновская длина волны, соответствующая \(m_\gamma\), равна
	\[
	\lambda_c = \frac{\hbar}{m_\gamma c} \approx 2.7\times10^{11}~\text{м},
	\]
	что на много порядков больше атомного облака (типичный размер \(\sim 1\) мм). Следовательно, эффекты распространения, обусловленные конечной массой фотона, в этом эксперименте полностью пренебрежимы. Отрицательная временная задержка является чисто координационным эффектом, независимым от абсолютной шкалы масс.
	
	Обратно, тот факт, что один фреймворк правильно предсказывает и:
	\begin{itemize}
		\item существование отрицательных групповых задержек с величиной, определяемой \(\ln\Keff\) (настоящее Приложение),
		\item верхний предел на массу фотона (\(m_\gamma < 10^{-18}\) эВ) и его дискретный узел \(N_\gamma = -410\) (Приложение Photon Mass),
	\end{itemize}
	представляет собой мощную проверку согласованности. Эти два предсказания относятся к совершенно разным физическим областям (квантовая оптика vs. свойства фундаментальных частиц), но разделяют одни и те же фундаментальные константы \(q\), \(S_{\mathrm{odd}}/S_{\mathrm{even}}\) и понятие эффективности координации \(\Keff\). Для их совместимости не требуется никакой подгонки.
	
	Таким образом, эксперимент с отрицательной задержкой фотонов не противоречит предсказанию YUCT массы фотона; напротив, оба служат независимыми опорами, поддерживающими единую координационную парадигму.
	
	\section{Отзыв предыдущего слепого предсказания}
	
	В более ранней версии этого Приложения (V.2.0) мы ошибочно предсказывали, что групповая задержка будет масштабироваться как \(T^{-0.20}\), что приводило к 13\%-ному изменению при удвоении температуры. Это предсказание основывалось на неверной линейной зависимости \(|\tau_g| \propto \Keff\). Правильное логарифмическое соотношение, выведенное здесь, показывает, что температурный эффект примерно на два порядка меньше. Поэтому утверждение о 13\%-ном эффекте недействительно и официально отзывается. Приносим извинения за это упущение и благодарим соавторов математического исправления за обнаружение несоответствия.
	
	Тем не менее, основная концепция YUCT – что отрицательные групповые задержки являются проявлением эффективности координации – остаётся справедливой. Экспериментальное подтверждение отрицательных времён атомного возбуждения уже поддерживает идею о том, что фотон и атом разделяют скоординированный словарь. Исправленный закон масштабирования теперь предоставляет уточнённый теоретический framework для будущих проверок.
	
	\section{Заключение}
	
	Мы исправили YUCT-анализ эксперимента с отрицательной задержкой фотонов. Универсальный закон ошибок YUCT является логарифмическим, и соответственно связь между групповой задержкой и эффективностью координации должна быть логарифмической: \(|\tau_g| = \alpha_0 \ln\Keff\). Это приводит к очень слабой температурной зависимости (логарифмической, с относительным изменением порядка нескольких процентов при удвоении температуры) и устраняет ранее заявленный 13\%-ный эффект. Ошибочное предсказание отзывается. Альтернативные проверки, такие как изменение оптической глубины или длительности импульса, предлагают гораздо более чувствительные зонды для логарифмического масштабирования. Существование отрицательных временных задержек остаётся поразительным подтверждением скоординированной природы квантовых взаимодействий, и YUCT предоставляет последовательную информационно-теоретическую интерпретацию.
	
	\section*{Благодарности}
	
	Автор благодарит соавторов математического исправления за указание на несоответствие между линейной гипотезой и логарифмической структурой YUCT. Также благодарит участников семинара YUCT за стимулирующие обсуждения и группу из Торонто (Aephraim Steinberg, Daniela Angulo и коллеги) за их прекрасный эксперимент.
	
	\section*{Доступность данных}
	
	Все вычисления и дополнительные материалы доступны по адресу \url{https://github.com/Alexey-Yakushev-YUCT/YPSDC}.
	
	\begin{thebibliography}{99}
		
		\bibitem{Angulo2024}
		D. Angulo, K. Thompson, A. Jiao, H. Wiseman, A. Steinberg, ``Experimental evidence that a photon can spend a negative amount of time in an atom cloud,'' arXiv:2409.03680 [quant-ph] (2024).
		
		\bibitem{Nixon2024}
		V.-M. Nixon et al., ``Measuring the time transmitted photons spend as atomic excitations,'' American Physical Society DAMOP Conference, Session Y07 (June 2024).
		
		\bibitem{Yakushev2026a}
		A. V. Yakushev, ``Yakushev Unified Coordination Theory (YUCT),'' Zenodo, \url{https://doi.org/10.5281/zenodo.18444599} (2026).
		
		\bibitem{Yakushev2026b}
		A. V. Yakushev, ``Appendix P: Temperature Dependence of Gravity,'' Zenodo (2026).
		
		\bibitem{Yakushev2026c}
		A. V. Yakushev, ``Appendix L: Fractal Error Scaling,'' Zenodo (2026).
		
		\bibitem{Yakushev2026d}
		A. V. Yakushev, ``Appendix X: Generalization of Shannon Information Theory,'' Zenodo (2026).
		
		\bibitem{Thompson2023}
		K. Thompson et al., ``How much time does a photon spend as an atomic excitation before being transmitted?'' arXiv:2310.00432 (2023).
		
		\bibitem{Brillouin1960}
		L. Brillouin, \textit{Wave Propagation and Group Velocity} (Academic Press, 1960).
		
		\bibitem{Milonni2004}
		P. Milonni, \textit{Fast Light, Slow Light and Left-Handed Light} (CRC Press, 2004).
		
		\bibitem{Wiseman2023}
		H. M. Wiseman, A. M. Steinberg, and M. Hallaji, ``Obtaining a single-photon weak value from experiments using a strong coherent state,'' AVS Quantum Sci. \textbf{5}, 024401 (2023).
		
	\end{thebibliography}
	
\end{document}